(1) 令入射角为 \(i\)、出射角为 \(e\)，棱镜内部两折射角为 \(r_1,r_2\)，则几何关系 \(r_1+r_2=A\)。总偏向是光线两次转向之和，\(\delta=i+e-A\)；界面折射满足 \(\sin i=n\sin r_1\)、\(\sin e=n\sin r_2\)。对于对称平面棱镜以及同一种外介质，光路反演对称性使最小偏向发生在两侧几何对称时。
(2) 极小偏向条件为 \(i=e\) 且 \(r_1=r_2=A/2\)。由 \(\delta_{\min}=2i-A\) 得 \(i=(A+\delta_{\min})/2\)，代入第一次折射的 Snell 定律便得 \(n=\sin[(A+\delta_{\min})/2]/\sin(A/2)\)。适用条件是光能从棱镜出射且最小偏向位置真实存在，例如 \(n\sin(A/2)\le1\)，否则形式解超出了物理角度范围。
(3) 由 \(\delta_{\min}(\lambda)=2\arcsin[n(\lambda)\sin(A/2)]-A\)，逐项求导为 \(d\delta_{\min}/d\lambda=2\sin(A/2)[dn/d\lambda]/\sqrt{1-n^2\sin^2(A/2)}\)。对正常色散材料 \(dn/d\lambda<0\)，波长越短偏向角越大。这个导数同时指出接近出射几何边界时形式角度色散变大，但透射效率与像差不一定允许无限提高分辨率。
(4) 两线相隔小波长 \(\Delta\lambda\) 时棱镜的角分离约为 \(|\Delta\delta|=|d\delta_{\min}/d\lambda|\,\Delta\lambda\)。若出射光束受直径 \(D\) 的近似圆形孔径限制，衍射首暗环角尺度为 \(1.22\lambda/D\)，瑞利式估计需 \(\Delta\lambda\gtrsim1.22\lambda/[D|d\delta_{\min}/d\lambda|]\)。这只是衍射受限估算，探测器采样和实际像差可能提高所需间隔。
